% ── results/gen_paradigm.tex ──────────────────────────────────────────────────
\subsection{Generation Paradigm Comparison: AR, NAR, and Masked Diffusion}
\label{sec:res_gen_paradigm}

\begin{table}[h]
\centering
\caption{Comparison of generation settings built on the same frozen JEPA encoder.
  AR v4 is the main 3-condition baseline.
  NAR and masked diffusion test new generation paradigms with the same
  3-condition target.
  AR v6 keeps the AR decoder but extends conditioning from 3 to 7 properties.
  Novelty = fraction of generated sequences with no exact match to the
  868k training corpus.}
\label{tab:gen_comparison}
\resizebox{\textwidth}{!}{%
\begin{tabular}{lccccc}
\toprule
\textbf{Paradigm} & \textbf{Charge $R^2$} & \textbf{Charge MAE}
  & \textbf{Length $R^2$} & \textbf{Novelty} & \textbf{Unique} & \textbf{Key note} \\
\midrule
AR Decoder (v4, proposed)    & $\mathbf{0.866}$ & $\mathbf{2.23}$ & $-1.71$ & 1.000 & 0.946 & best charge control \\
NAR Decoder                  & 0.693 & 3.566 & $-1.155$ & 0.954 & 0.727 & weaker than AR on charge and length \\
Masked Diffusion             & 0.702 & 2.525 & $\mathbf{1.000}$ & 0.955 & 0.929 & best new paradigm; exact length control \\
AR v6 (7-dim conditions)     & 0.659 & 1.630 & $-0.605$ & 1.000 & 0.925 & adds partial helix/pI control \\
\bottomrule
\end{tabular}
}
\end{table}

All three decoders share the same frozen JEPA encoder (Section~\ref{sec:pretrain})
and use matched evaluation targets.
The AR decoder (v4) remains the strongest model for charge control
($R^2{=}0.866$).
NAR keeps reasonable charge control ($R^2{=}0.693$), but it is clearly weaker
than AR and also weak on length control.
Masked diffusion gives the best result among the new paradigms.
Its charge control is lower than AR but still strong ($R^2{=}0.702$), and it
matches target length exactly in this evaluation ($R^2{=}1.000$).

The v6 model answers a different question: what happens when we add four new
control dimensions on top of the AR decoder.
It keeps moderate charge control ($R^2{=}0.659$) and shows partial control of
helix propensity ($R^2{=}0.533$) and pI ($R^2{=}0.330$).
However, it does not control GRAVY, hydrophobic moment, or length well.
The AMP-score signal is present but weak ($R^2{=}0.237$).

\finding{The generation paradigm matters.
The original AR decoder is still best for charge control, but masked diffusion
is the strongest new paradigm because it keeps good charge control and gives
perfect length control in this test.
The v6 model shows that adding more condition dimensions does not guarantee
reliable control of every property.}
